\documentclass[11pt]{article}
\usepackage{fontspec}
\setmainfont{Times New Roman}
\setsansfont{Arial}
\setmonofont{Consolas}
\usepackage{amsmath,amssymb}
\usepackage{geometry}
\usepackage{microtype}
\geometry{a4paper,margin=3cm}
\setlength{\parindent}{1.25em}
\setlength{\parskip}{0pt}
\emergencystretch=6em
\allowdisplaybreaks

\begin{document}

\begin{center}
{\Large\bfseries 4. К инвариантној теорији форм в $n$ пременных}\par
\vspace{1.2em}
\emph{Journal f. d. reine u. angew. Math.} 139 (1911), с. 118--154
\end{center}

\section*{Ввод.}
В пројективној инвариантној теорији форм в $n$ пременных пытаньа, присојединьене к главному пытаньу -- доказу конечности система форм, -- мало былы размотрене, ако идемо поза бинарну и тернарну област. То сут пытаньа о взаимној зависимости форм и о методах овладаньа тым связом меджу формами. Те методы основно опирајут се -- како то в бинарној и тернарној области јест полно доказано -- на две теоремы, кторе Студы означил како ``фундаменталне теоремы симболичного метода'', имено на теорему о симболичној представимости форм и на теорему о симболичных идентичностјах.\footnote{Е. Студы, \emph{Methoden zur Theorie der ternären Formen} (Леипзиг, Teubner, 1889). II. \S\ 6. -- О наступаньу тыцх фундаменталных теорем такоже в инвариантној теорији специалных трансформацијских груп срав. Студы, \emph{Das Apollonische Problem} (\emph{Math. Ann.} Бд. 49 [1897]) и \emph{Zur Differentialgeometrie der analytischen Kurven} (\emph{Transactions of the American Math. Society}, Бд. 1 [1909]).} Последньа теорема говори, же все релације меджу целыми инвариантными творјеньами достајут се последовательным применьеньем малого числа симболичных идентичностиј; значит, симболичны метод, и само он, даје познанье связа меджу формами без выхода из области инвариантов.

На причину трудности при преходу к $n$ пременным уже указал Gordan в својем Ерлангенском програму;\footnote{П. Gordan, \emph{Über das Formensystem binärer Formen}. С. 46. (Леипзиг, Teubner, 1875).} она лежи в том, же прва теорема симболичного метода в својем обчем формулованьу вече не плати. В $n$-арној области, како познано, наставајут ряды пременных вышего нивоја $p_\rho$\footnote{За означенье срав. \S\ 1.} и аналогично ряды симболов вышего нивоја $S^\rho$; к тым рядам се доходит и тогда, когда се, како изворно при Цлебсцху,\footnote{А. Clebsch, \emph{über eine Fundamentalaufgabe der Invariantentheorie}. (\emph{Abh. der Ges. d. Wiss. zu Göttingen}, Бд. XVII, 1872).} изходит из форм с либоволно многыми рядами $p_\rho$, и тогда, когда се, по поступку Ф. Деруытса и А. Цапеллија,\footnote{Срав. А. Цапелли, \emph{Lezioni sulla teoria delle forme algebriche} (Napoli, 1902), \S\S\ 22--24, и там наведену литературу.} изходит из форм, кторе обсахујут само либоволно много когредиентных рядов $x$. Симболично представјенье, кторе бы експлицитно\footnote{Точнејше определьенье ``експлицитного представјеньа'' в \S\ 2.} обсаховало ряды $p_\rho$, $S^\rho$, тепер не обстаје; потребно јест разложити $p_\rho$, $S^\rho$ в поједине ряды $x$ и в к ним контрагредиентне $s$ и разделити е меджу различне детерминанты. Правда, с помочју диференциалных условиј (срав. формулу (19.)) можно проверити, же даны агрегат детерминантов има властност зависети само од рядов $p_\rho$, $S^\rho$; але обзор над целоју множиноју инвариантных творјень и над процесами за их творјенье не можно таким способом добити.\footnote{Такоже рачунско велми ценно дальне развитје симболикы, кторе изходит од Студыја (сообчено од Ф. Енгела, \emph{Ber. d. K. Sächs. Ges. d. Wiss., math.-phys. Kl.}, 1900, С. 126, и за кватернарну област од Е. Студы, \emph{Geometrie der Dynamen}, С. 135), не јест експлицитно представјенье в нашем смыслу. Оно бы, на примјер, едва привело к несимболичным процесам диференцираньа за творјенье форм.}

Тепер показујемо в \S\S\ 2 и 6, како с применьеньем ``теоремы о одговарјајучих матрицах''\footnote{Ако изоставити Грассманнову \emph{Ausdehnungslehre} из 1862, Нр. 112, најпрво при А. Бриллу, \emph{über zwei Berührungsprobleme}, \S\ 2 (\emph{Math. Ann.} Бд. 4, 1871).} можно знову добити прву фундаменталну теорему в ей обчности, заменивши представјенье чрез детерминанты представјеньем чрез ``матричне продукты''. В \S\ 2 показује се, же всак матричны продукт, кторы експлицитно обсахује ряды $S^\rho,p_\rho$, јест инвариантно творјенье основној формы; обратност, дана в \S\ 6, же всако инвариантно творјенье можно експлицитно представити матричными продуктами, успева само чрез пренос другој фундаменталној теоремы (\S\S\ 3--5).

Та теорема јест познана за формы, кторе обсахујут само ряды пременных $x$.\footnote{Е. Пасцал, \emph{Giorn. di mat.} 26 (1888), С. 33, 102; \emph{Rendic. d. Accad. d. Linc.} (4) 4 (1888), С. 119; иб. \emph{Mem.} (4) 5 (1888), С. 375.} Обче пытанье -- установити целу множину неразложимых идентичностиј, при кторых все ряды $S^\rho,p_\rho$ разматривајут се како неразложиме, -- формуловал Студы;\footnote{\emph{Methoden \dots}, С. 204, заметка 17).} он едновременно види в числу наставајучих идентичностиј меру трудности методов в $n$-арној области. Когда тепер успева се сојединити целу множину идентичностиј в једну формулу (36.), та трудност јест најманье снижена до минимума. При применьеньу вшак често наставајут некторе вызначене специалне случаји од (36.), како (20.), (21.), кторе сут карактеризоване тым, же допушчајут експлицитно представјенье без субституције нових рядов; або формула (31.), ктора называ се ``разкладна идентичност'', бо ряд $p_\rho$ в ньеј на вид јавја се разложеным в ряды $y$.

На двех фундаменталных теоремах можно тепер лахко изградити дальну теорију. Прва теорема непосредне даје процесы творјеньа форм, симболичны процес свијаньа и несимболичне процесы диференцираньа,\footnote{В работе \emph{Invariante Gebilde von Nullsystemen} (\emph{Sitzungsber. d. Ak. d. Wiss. Wien}, Бд. 97, Абт. ИИа, 1888) Мертенс е другим, не непосредне обчим способом дошел до кватернарного процеса свијаньа. Там наставајучи алтернирајучи инвариант од 6 рядов $S^2$ разликује се од нашего, кторы наставаје при једнократном применьеньу субституције (11.), чрез агрегат парных инвариантов. За једен специалны случай кватернарны процес свијаньа најде се такоже при П. Гордану, \emph{Das simultane System von zwei quadratischen quaternären Formen} (\emph{Math. Ann.} Бд. 56, 1903).} меджутым како разкладна идентичност меджу тыми процесами изкльучаје все еквивалентне (\S\ 6). Познанье процесов диференцираньа далеј даје пренос појма ``нормална форма'' на $n$-арну област и, с помочју доказного поступка, даного Мертенсом\footnote{Ф. Мертенс, \emph{Über invariante Gebilde quaternärer Formen}. (\emph{Sitzber. d. Ak. d. Wiss. Wien}. Бд. 98, Абт. ИИа, 1889.)} в кватернарној области, теорему, же систем инвариантных творјень можно заменити еквивалентным системом нормалных форм (\S\ 7). Под тым последньим предположеньем ја показала в мојеј дисертацији,\footnote{\emph{Über die Bildung des Formensystems der ternären biquadratischen Form}. (Тутој журнал, Бд. 134, 1908.)} же тернарны процес свијаньа можно створити последовательным применьеньем двех основных свијань. На основе тој теоремы и теорије разложень в ряды ја потом, введши појем ``ряд форм'', развила главне редукцијне теоремы за системы форм. Целком аналогично в $n$-арној области процес свијаньа можно створити из $n-1$ основных свијань (\S\ 8), и можно установити редукцијне теоремы (\S\ 9).

\paragraph{Додатна заметка.} При Салзбургском сообченьу проф. \mbox{E. Müller} из Ведньа обратил моју внимльивост на то, же рачунанье с матричными продуктами, кторе лежи в основе тој работы, в принципу јест идентично с калкулом Грассманновеј наукы о разширјеньу.\footnote{\emph{Hermann Graßmanns gesammelte mathematische und physikalische Werke. Ersten Bandes zweiter Teil: Die Ausdehnungslehre von 1862. In Gemeinschaft mit H. Graßmann d. Jüngeren herausgegeben v. Fr. Engel} (Леипзиг, Teubner, 1896).} В истине Грассманнов ``комбинаторичны продукт'' $[S^{\rho_1}S^{\rho_2}\cdots S^{\rho_i}]$ можно разматривати како матрицу, дану тыми рядами (срав. \S\ 2), и то ``прогресивны продукт'' како матрицу самых рядов, ``регресивны'' како матрицу прирядженых рядов $S_{n-\rho_1},S_{n-\rho_2},\ldots,S_{n-\rho_i}$, меджутым како матрица, одговарјајуча ``смешаному продукту'', наставаје, когда неколико рядов $S$ сојединьа се чрез субституцију (9.) в нове ряды $P,Q$. Наш ``приряджены ряд'' одговарја при том Грассманновому ``дополньеньу'', а наш ``матричны продукт'' -- ``смешаному продукту ступньа 0''.

По тој приряджености разложенье, дано в \emph{Ausdehnungslehre} 1862, Нр. 113, ставаје идентично с формулоју, дуалистичној к нашој формули (32.). В специалном случају $\rho=1$ формула (32.) преходи в (17.), а дуалистична формула в (16.). Из тыцх специалных случајев Грассманнового разложеньа \mbox{Müller} недавно извел выше споменуте идентичности (20.), (21.).\footnote{Е. \mbox{Müller}, \emph{Beiträge zur Graßmannschen Ausdehnungslehre. 1. Mitteilung: Einige allgemeine Sätze}. (\emph{Sitzungsber. d. Ak. d. Wiss. Wien}; Бд. 118, Абт. ИИа, Јули 1909), Сатз 16 и 18.}

Наше изводженье, в разлике од Грассманн--Муллерового, едновременно даје доказ полности идентичностиј, што за пытаньа, кторе туто сут разматриване, јавја се најважнејшим; оно иде од (20.), (21.) далеј к најобчејшим неразложимым идентичностјам (36.), кторе до сега часа, како се зда, еще не былы замечене.

\end{document}

